\section{Validity Statement}

The \vm{} and Kubernetes environments are not identical. They differ in scheduling
mechanisms, networking, and resource allocation. The goal is not to eliminate these
environmental differences but to measure their combined effect as part of real system
behaviour, which is how architectural migration decisions are made in practice.

Running both environments on the same physical host removes hardware variance as a
confounding factor. The fixed Kind single-node setup removes autoscaling variance, ensuring
that observed differences arise from the execution model (Docker container orchestration on
a single-host daemon vs.\ Kubernetes pod scheduling with per-pod resource limits) rather
than from dynamic scaling events.

All experiments use \dagster{} as the representative workflow orchestration system and
synthetic workloads. Results may not generalise directly to all workflow systems or all
infrastructure configurations. However, the behavioural patterns identified --- the failure
degradation curve, the blast radius reduction, and the crossover point --- are expected to
be directionally relevant for similar pipeline orchestration frameworks under comparable
concurrency conditions.

If results are inconsistent across experimental runs, that inconsistency is documented and
analysed as a finding rather than suppressed as error.
